\section{基于波动力学的类别偏置预测模型设计}
\subsection{模型整体框架}
本文模型由特征提取主干网络、类别相似度邻接矩阵构建模块、波动演化预测模块、动态损失补偿模块四部分组成。
主干网络负责从输入图像中提取各类别特征；邻接矩阵用来刻画类别之间特征相似度，作为波动传播的空间连接关系；波动演化模块以训练迭代为时间步，模拟类别偏置信号沿着类别关联关系传播、累积、衰减，预测未来时刻各类别的偏置量；最后根据预测得到的偏置值动态调整各类损失权重，抑制偏置持续放大。

\subsection{类别相似度邻接矩阵构建}
设一共存在 $C$ 个类别，每一轮训练结束后，统计每个类别在特征空间的中心向量。类别 $i$ 和类别 $j$ 的相似度采用余弦相似度计算：
\[
A_{ij}=\frac{\boldsymbol{\mu}_i^\top \boldsymbol{\mu}_j}{\|\boldsymbol{\mu}_i\|\|\boldsymbol{\mu}_j\|}
\]
将相似度矩阵归一化后得到邻接矩阵 $\boldsymbol{A}$，代表偏置信号可以在类别之间传播的连接强度。相似度越高，两类之间偏置相互影响越强。

\subsection{偏置波动演化方程}
将每一次训练迭代看作一个时间步 $t$，$\boldsymbol{b}_t\in\mathbb{R}^C$ 代表 $t$ 时刻各类别的偏置向量。借鉴波动扩散思想，离散演化公式设计为：
\[
\boldsymbol{b}_{t+1} = \alpha \boldsymbol{b}_t + \beta \boldsymbol{A}\boldsymbol{b}_t + \gamma \boldsymbol{\epsilon}_t
\]
- $\alpha$：偏置自身保留系数；
- $\beta$：跨类别传播系数；
- $\gamma$：噪声衰减系数；
- $\boldsymbol{\epsilon}_t$：本轮新产生的原始偏置。

通过迭代推演，可以预测后续若干步的偏置变化趋势，提前获知哪些类别偏置会快速增大。

\subsection{动态损失补偿策略}
传统加权损失权重固定不变，无法跟随偏置演化动态调整。本文利用预测偏置 $\hat{b}_t^i$ 计算类别 $i$ 的动态权重：
\[
w_t^i = \exp(\tau\cdot \hat{b}_t^i)
\]
$\tau$ 为控制补偿强度的超参。偏置预测值越大，该类别损失权重越高，迫使模型在后续训练中重点修正该类别的偏差。最终损失函数：
\[
\mathcal{L}=-\sum_{i=1}^C w_t^i \cdot y_i \log(p_i)
\]

\subsection{训练流程}
1. 前向传播，计算各类特征中心、当前偏置；
2. 更新类别邻接相似度矩阵；
3. 代入波动演化方程，预测下一时刻偏置；
4. 根据预测偏置更新各类损失权重；
5. 反向传播更新主干网络参数；
6. 重复迭代直到收敛。
